\section{Related Work}
\label{sec:related}

\noindent\textbf{Clinical Gait Analysis.}
Observational gait analysis has long been the cornerstone of clinical assessment for children with CP~\cite{gage2009treatment}. The Edinburgh Visual Gait Score (EVGS)~\cite{read2003edinburgh} provides a structured 17-item scoring system using three-point ordinal scales, enabling systematic evaluation from sagittal and coronal video recordings. Winters~\etal~\cite{winters1987gait} established the foundational classification of hemiplegic gait into four types, later extended to bilateral patterns by Rodda and Graham~\cite{rodda2004classification}, identifying True Equinus, Jump Gait, Apparent Equinus, and Crouch Gait as canonical subtypes. Despite their clinical utility, these classification systems exhibit limited inter-rater reliability~\cite{wren2005reliability}, motivating automated alternatives.

\noindent\textbf{Skeleton-Based Action Recognition.}
Human pose estimation has advanced rapidly with deep learning approaches. OpenPose~\cite{cao2017realtime} and HRNet~\cite{sun2019deep} established accurate 2D pose estimation, while recent foundation models such as Sapiens~\cite{khirodkar2024sapiens} achieve state-of-the-art keypoint detection with 133 COCO-WholeBody~\cite{jin2020whole} keypoints. For skeleton-based action recognition, Spatial-Temporal Graph Convolutional Networks (ST-GCN)~\cite{yan2018spatial} model skeleton sequences as spatiotemporal graphs, with subsequent improvements through attention mechanisms~\cite{shi2019two} and transformer architectures~\cite{plizzari2021skeleton, 11484350}. However, these deep learning approaches typically require large-scale training data, making them challenging to apply directly to small clinical datasets.

\noindent\textbf{Automated Gait Analysis.}
Computational approaches to gait analysis have evolved from instrumented walkways and marker-based motion capture to markerless video-based methods~\cite{stenum2021two, kidzinski2020deep}. Kidzinski~\etal~\cite{kidzinski2020deep} demonstrated that deep learning can predict clinical gait metrics from 2D video, while Stenum~\etal~\cite{stenum2021two} validated markerless pose estimation for clinical gait analysis. For pediatric populations specifically, several works have explored automated CP gait classification~\cite{zhang2023automated, chen2024gait}, though most rely on instrumented gait laboratories rather than standard video. The CV4CHL challenge~\cite{cv4chl2026} uniquely addresses this gap by providing smartphone-captured multi-view keypoint data in a clinical setting.

\noindent\textbf{Ensemble Learning and Domain Knowledge Integration.}
Ensemble methods combining multiple base learners---including Gradient Boosting~\cite{friedman2001greedy}, Random Forests~\cite{breiman2001random}, and Extra Trees~\cite{geurts2006extremely}---have proven effective for tabular data with limited samples. The integration of domain knowledge into machine learning pipelines has been explored through physics-informed neural networks~\cite{raissi2019physics}, clinical rule injection~\cite{rudin2019stop}, and hybrid expert-ML systems~\cite{rajkomar2019machine}. Our work extends this direction by demonstrating that rule-based clinical knowledge can be effectively integrated at the prediction layer rather than the feature layer, a distinction that proves critical for small-sample regimes.


\noindent\textbf{Adjacent Multimodal Learning and Supervision.}
Research on resource-constrained AI and multimodal learning provides broader context for practical visual-analysis systems. Energy-efficient retrieval-augmented generation with small language models has been surveyed for on-device settings~\cite{cheng2026toward}. In other application domains, in-prompt process supervision for short-video moderation~\cite{liu-etal-2026-ips} and reasoning-enhanced multimodal domain-adaptive pretraining for video content governance~\cite{wang-etal-2025-reasoning-enhanced} investigate supervision and representation adaptation. Audio-enhanced vision--language modeling with latent-space broadening has also been studied for data expansion~\cite{10.1145/3711896.3737195}. These works do not address pediatric gait assessment or establish the clinical-rule fusion method used here; they are cited only as adjacent methodological context.
